\documentclass[11pt,a4paper]{article}

% ── Packages ──
\usepackage[utf8]{inputenc}
\usepackage[T1]{fontenc}
\usepackage{amsmath,amssymb}
\usepackage{graphicx}
\usepackage{booktabs}
\usepackage{hyperref}
\usepackage{geometry}
\usepackage{caption}
\usepackage{float}
\usepackage{enumitem}
\usepackage{array}

\geometry{margin=2.5cm}
\hypersetup{colorlinks=true, linkcolor=blue, citecolor=blue, urlcolor=blue}

% ── Title ──
\title{%
  Chaos-Informed NGRC--TCN Encoder--Decoder\\
  for Multi-Step Waveform Forecasting\\[0.3em]
  \large Validation on the Lorenz System\\[0.5em]
  \normalsize NU Research Internship 2026 --- Time Series Forecasting Group
}
\author{
  Ibrahim Hanafy\\
  Zewail City of Science and Technology
}
\date{August 2026}

\begin{document}
\maketitle

% ════════════════════════════════════════════════════════════════
\begin{abstract}
This report documents the design, diagnosis, and successful validation of a chaos-informed forecasting architecture on the Lorenz system.
The model applies Phase Space Reconstruction (PSR) and a Next-Generation Reservoir Computing (NGRC) polynomial feature map as a zero-parameter front-end to a Temporal Convolutional Network (TCN) encoder--decoder.
After correcting an input-length pathology in the decoder, the model reached $R^2 = 0.995$ at 0.23~Lyapunov times and $R^2 = 0.956$ at one full Lyapunov time, outperforming a raw-input TCN control (0.738), a recursive NGRC ridge model (0.157), a linear autoregressive baseline (0.246), and persistence ($-0.434$).
Validation includes a shuffled-target leakage control, attractor reconstruction, and Lyapunov exponent recovery.
\end{abstract}

\tableofcontents
\newpage

% ════════════════════════════════════════════════════════════════
\section{Introduction}

Chaotic time series prediction has become one of the most actively studied problems in nonlinear dynamics and machine learning.
Many systems in nature --- fluid flows, weather, cardiac rhythms, financial markets --- exhibit sensitive dependence on initial conditions, making long-horizon forecasting fundamentally difficult.
Quantifying predictability in terms of Lyapunov time provides a principled, system-independent measure of forecast difficulty.

This report is organised around an experimental narrative rather than a single result.
Section~2 reviews the relevant literature.
Section~3 states the architecture.
Section~4 documents every trial in chronological order, including the failures, because the diagnostic path is the main contribution.
Section~5 states the findings and Section~6 the limitations.

\subsection{Contributions}

\begin{enumerate}[leftmargin=*]
  \item An NGRC polynomial front-end feeding a MIMO TCN encoder--decoder, evaluated against a matched raw-input control rather than against unrelated baselines.
  \item Evaluation across a sweep of forecast horizons expressed in Lyapunov time, rather than at a single arbitrary horizon.
  \item Identification of an input-length pathology in interpolation-based decoders: reducing input length from 540 to 64 samples raised $R^2$ at one Lyapunov time from 0.34 to 0.96.
\end{enumerate}


% ════════════════════════════════════════════════════════════════
\section{Literature Review}

\subsection{Chaotic Time Series and Phase Space Reconstruction}

Chaotic signals arise from nonlinear deterministic systems and are characterised through Lyapunov exponents, correlation dimension, and attractor geometry~\cite{ramadevi2022}.
Takens' delay embedding theorem provides the formal basis for reconstructing an attractor from a scalar observable:
\begin{equation}
  \mathbf{v}(t) = \bigl[x(t),\; x(t-\tau),\; \ldots,\; x(t-(d-1)\tau)\bigr],
\end{equation}
where the delay $\tau$ is conventionally selected from the first minimum of Average Mutual Information (AMI) and the embedding dimension $d$ from the False Nearest Neighbours (FNN) criterion.

Phien et al.~\cite{phien2023} applied AMI and FNN before feeding embedded vectors to LSTM autoencoders, and compared five multi-step strategies (recursive, direct, MIMO, DirRec, DirMO), concluding that MIMO performed best and recursive worst due to error accumulation.
This directly motivates the MIMO design adopted here.

\subsection{Reservoir Computing and NGRC}

Reservoir computing fixes a high-dimensional nonlinear dynamical system and trains only a linear readout.
Gauthier et al.~\cite{gauthier2021} removed the random reservoir entirely, constructing features directly from delayed observations and their polynomial cross-terms.
On Lorenz-63 this achieved NRMSE $\approx 1.06 \times 10^{-4}$ with training in under 10\,ms.
Critically, NGRC operates recursively: a one-step map is fitted, then the prediction is fed back so the system runs autonomously.

Shahi et al.~\cite{shahi2022} formalised NGRC as a Nonlinear Vector Autoregression and compared LSTM, GRU, ESN and NVAR on chaotic systems including zebrafish cardiac voltage recordings, finding gated RNNs to underperform on chaotic forecasting relative to reservoir approaches.

An important limitation for the univariate case is documented in the partial-observation NGRC literature: when only one state variable is observed and time-delay coordinates are used, degree-2 monomials are insufficient, and higher-degree bases (e.g.\ Chebyshev polynomials of degree~9) are required.
This predicts that a degree-2 univariate NGRC will underperform, which is confirmed experimentally in Section~4.

\subsection{Temporal Convolutional Networks}

TCNs offer parallelisable training, stable gradients, and exponentially growing receptive fields via dilated causal convolutions.
Dudukcu et al.~\cite{dudukcu2023} benchmarked hybrid TCN--LSTM and TCN--GRU models against twelve alternatives on Lorenz, R\"ossler, and a Lorenz-like system, reporting average RMSE 0.0022 on chaotic data.
Their prediction horizon is one time step from a 10-step history, which on the Lorenz system at $dt = 0.01$ corresponds to 0.009~Lyapunov times --- well within the trivially predictable regime.

Fu et al.~\cite{fu2024} proposed FGNet, a Fourier-attention encoder with a GRU decoder, evaluated at 1, 4, 7 and 10 steps on Mackey--Glass, Lorenz and sunspot series.
Their reported Lorenz $R^2$ values are higher than those obtained here at matched horizons (Section~4.4).
However, their reported Lorenz largest Lyapunov exponent is 0.012285, whereas the accepted value for $(\sigma, \rho, \beta) = (10, 28, 8/3)$ is 0.9056 --- a factor of roughly 74, consistent with a per-sample versus per-time-unit unit convention that is not stated.
With the correct exponent, their maximum horizon of 10 steps at $dt = 0.01$ corresponds to 0.09~Lyapunov times, i.e.\ entirely within the short-term-predictable regime.

\subsection{Identified Gap}

Across the reviewed work, three patterns recur: evaluation at a single short horizon, absence of naive baselines, and horizons that are not expressed in units comparable across systems.
No reviewed work combines an NGRC polynomial front-end with a deep TCN encoder--decoder for multi-step waveform forecasting, nor reports a predictability curve in Lyapunov time.
That is the gap addressed here.


% ════════════════════════════════════════════════════════════════
\section{Method}

\subsection{Architecture}

The pipeline is:
\[
\text{Signal} \;\to\; \text{PSR} \;\to\; \text{NGRC feature map} \;\to\; \text{TCN encoder} \;\to\; \text{interpolate} \;\to\; \text{TCN decoder} \;\to\; \text{forecast.}
\]

\textbf{Chaos front-end (zero learnable parameters).}
The delay embedding is expanded with all unique degree-2 monomials:
\begin{equation}
  \Phi(t) = \bigl[\underbrace{v_1, \ldots, v_d}_{d\;\text{linear}},\; \underbrace{v_1^2, v_1 v_2, \ldots, v_d^2}_{d(d+1)/2\;\text{quadratic}}\bigr].
\end{equation}
The quadratic terms are motivated by the structure of the Lorenz equations themselves, which contain the products $xy$ and $xz$; a purely linear map cannot represent these.

\textbf{Encoder.} Four residual TCN blocks, dilations $[1, 2, 4, 8]$, kernel~7, 64 channels, causal.

\textbf{Decoder.} Four non-causal blocks of the same structure.
The encoder output is resampled along the time axis by linear interpolation to the target length, followed by a $1 \times 1$ convolution to a single channel.
This interpolation is the source of the pathology reported in Section~4.2.

\subsection{Evaluation Protocol}

Every experiment reports at minimum: the model, a zeros/mean baseline, and a persistence baseline.
Horizons are converted to Lyapunov time via $\lambda_{\max} t$.
On Lorenz, $\lambda_{\max} = 0.9056$ and $dt = 0.01$, giving one Lyapunov time $= 110.4$ samples.


% ════════════════════════════════════════════════════════════════
\section{Experimental Trials}

\subsection{Trial 1 --- Initial Lorenz Validation}

The architecture was applied to the univariate $x(t)$ component of the Lorenz system: deterministic, noise-free, with a textbook $\lambda_{\max}$.

PSR parameters converged cleanly ($\tau = 16$, $d = 3$, 9 features).

\begin{table}[H]
\centering
\caption{Trial~1, input length 540 samples.}
\begin{tabular}{rr|c|c|c}
\toprule
$H$ & $\lambda_{\max} t$ & Chaos $R^2$ & NGRC $R^2$ & Persistence $R^2$ \\
\midrule
25  & 0.23 & $+0.9602$ & $+0.6932$ & $+0.5345$ \\
110 & 1.00 & $+0.3633$ & $+0.1565$ & $-0.3092$ \\
550 & 4.98 & $-0.0135$ & $+0.0270$ & $-1.1075$ \\
\bottomrule
\end{tabular}
\end{table}

The architecture worked, but degraded sharply by one Lyapunov time.


\subsection{Trial 2 --- Input-Length Ablation}

A linear autoregressive control on 64 raw lags achieved $R^2 = 0.2565$ at $H = 110$, close to the model's 0.3633, prompting an ablation on input length at fixed $H = 110$.

\begin{table}[H]
\centering
\caption{Trial~2. Shorter input is dramatically better --- the opposite of the usual expectation that more context helps.}
\begin{tabular}{rc}
\toprule
\textbf{Input length} & $R^2$ \textbf{at} $H = 110$ \\
\midrule
64  & 0.9590 \\
128 & 0.9271 \\
200 & 0.8180 \\
320 & 0.4056 \\
540 & 0.3367 \\
\bottomrule
\end{tabular}
\end{table}

\textbf{Mechanism.}
The decoder resamples the encoder output from length $L$ to length $H$ by linear interpolation.
Output index $t$ is therefore anchored to input index $t \cdot L/H$.
With $L = 540$ and $H = 110$, the first predicted sample is constructed from the oldest position in the input window, 5.4 time units ($\approx 4.9$ Lyapunov times) in the past, where the signal is decorrelated from the present.
With $L = 64$ the entire window spans 0.64 time units and remains correlated.
The long window was not neutral padding; it actively misaligned the decoder.

A secondary consideration: with dilations $[1, 2, 4, 8]$ and kernel~7, the encoder receptive field is approximately $1 + 2(7-1)(1+2+4+8) = 181$ samples, so inputs beyond that length are partially unreachable regardless.


\subsection{Trial 3 --- Corrected Lorenz Sweep and Full Ablation}

With input length set to 64, the sweep was repeated with all arms.
Table~\ref{tab:trial3} reports $R^2$, RMSE, and MAE for each model.

\begin{table}[H]
\centering
\caption{Trial~3, input length 64. Full comparison with RMSE and MAE.}
\label{tab:trial3}
\begin{tabular}{rr|rrr|rrr|rrr|c|c}
\toprule
 & & \multicolumn{3}{c|}{\textbf{ChaosForecaster}} & \multicolumn{3}{c|}{\textbf{VanillaTCN}} & \multicolumn{3}{c|}{\textbf{Pure NGRC}} & \textbf{Lin.\ AR} & \textbf{Persist.} \\
$H$ & $\lambda t$ & $R^2$ & RMSE & MAE & $R^2$ & RMSE & MAE & $R^2$ & RMSE & MAE & $R^2$ & $R^2$ \\
\midrule
25  & 0.23 & 0.995 & 0.069 & 0.051 & \textbf{0.996} & \textbf{0.065} & \textbf{0.048} & 0.693 & 0.556 & 0.407 & 0.916 & 0.533 \\
110 & 1.00 & \textbf{0.956} & \textbf{0.210} & \textbf{0.130} & 0.738 & 0.512 & 0.287 & 0.157 & 0.919 & 0.702 & 0.246 & $-0.434$ \\
550 & 4.98 & \textbf{0.075} & \textbf{0.963} & \textbf{0.766} & 0.018 & 0.992 & 0.797 & 0.027 & 0.987 & 0.789 & 0.039 & $-1.108$ \\
\bottomrule
\end{tabular}
\end{table}

Three readings follow directly:

\begin{itemize}[leftmargin=*]
  \item \textbf{At 0.23 Lyapunov times the NGRC front-end contributes nothing.}
    $R^2 = 0.995$ versus 0.996 is within run-to-run variation ($\pm 0.003$ observed across repeated runs).
    Short-horizon Lorenz-$x$ is close to linearly predictable (0.916 for linear AR).
    VanillaTCN achieves the lowest RMSE (0.065) and MAE (0.048).

  \item \textbf{At one Lyapunov time the front-end contributes substantially.}
    $R^2 = 0.956$ versus 0.738 corresponds to an RMSE reduction from 0.512 to 0.210 (59\%).
    This is where the dynamics must actually be propagated forward, and where the explicit quadratic terms match the quadratic structure of the governing equations.

  \item \textbf{At five Lyapunov times all methods collapse.}
    Values near zero are not meaningfully distinguishable.
    This horizon was additionally underpowered (27 validation and 54 test windows) and did not converge.
\end{itemize}

The recursive NGRC ridge model, with only 10 coefficients, reached $R^2 = 0.693$ at $0.23\lambda_{\max}t$.
Its underperformance relative to the deep model is consistent with the documented limitation of degree-2 monomials under univariate partial observation.


\subsection{Validation Checks}

\textbf{Shuffled-target control.}
Training the same model on inputs paired with permuted targets yielded $R^2 = -0.0051$ against $+0.9955$ for the real model, confirming no information leakage through window indexing, the PSR offset, or normalisation.

\textbf{Attractor reconstruction.}
Stitched predictions embedded at lag $\tau = 16$ reproduce the two-lobe Lorenz structure, confirming the model captured attractor geometry rather than only pointwise amplitude.
The predicted embedding is visibly noisier at fine scale than the ground truth.

\textbf{Lyapunov exponent recovery.}
A Rosenstein estimator was first calibrated on the true series, recovering $\lambda_{\max} = 0.9105$ against the textbook 0.9056 (0.5\% error) with a fit window of $[50{:}200]$, stable across four $(d, \tau)$ settings.
Earlier fit windows returned values up to 1.62, an overestimate of 78\% --- a sensitivity that plausibly explains the divergent exponents reported in the literature.
Applied to predictions, the estimator returned 0.6257 against 0.9552 from the truth, a 34.5\% underestimate.
Two causes apply: squared-error training damps fine-scale divergence, and the stitched trajectory is a concatenation of teacher-forced segments rather than a free-running trajectory, which inflates the initial neighbour distance.

\textbf{Forecast horizon.}
Using the $\sigma$-crossing criterion (first step at which absolute error exceeds one standard deviation), the model tracked for a mean of 106 steps at $H = 110$, i.e.\ $\geq 0.96$ Lyapunov times.
This is a lower bound: 94\% of windows never diverged within the available horizon.


\subsection{Comparison at FGNet Horizons}

Re-running at the horizons used by Fu et al.~\cite{fu2024} gives a directly comparable setting (same system, same $dt$, scale-invariant metric):

\begin{table}[H]
\centering
\caption{FGNet achieves higher accuracy at these horizons.}
\begin{tabular}{rr|cc|c}
\toprule
$H$ & $\lambda_{\max} t$ & This work & FGNet & Persistence \\
\midrule
1  & 0.01 & 0.9965 & 0.999994 & 0.9970 \\
7  & 0.06 & 0.9947 & 0.999569 & 0.9421 \\
10 & 0.09 & 0.9957 & 0.998631 & 0.8923 \\
\bottomrule
\end{tabular}
\end{table}

Two observations.
First, at $H = 1$ persistence (0.9970) exceeds this work (0.9965), indicating the regime is close to trivial.
Second, the present model's $R^2$ is essentially flat across $H = 1, 7, 10$ while persistence falls from 0.997 to 0.892; this indicates a fixed error floor near RMSE $= 0.06$ arising from the decoder rather than from horizon difficulty, and represents an identified implementation gap.
All FGNet horizons fall below 0.1~Lyapunov times.


% ════════════════════════════════════════════════════════════════
\section{Findings}

\begin{enumerate}[leftmargin=*]
  \item \textbf{The NGRC front-end provides substantial benefit at the one-Lyapunov-time horizon.}
    $R^2 = 0.956$ versus 0.738 for the matched raw-input control, a 59\% RMSE reduction.
    At shorter horizons (0.23 Lyapunov times) the benefit vanishes, consistent with the signal being near-linearly predictable.

  \item \textbf{Input-length pathology in interpolation-based decoders.}
    Reducing input from 540 to 64 samples raised $R^2$ at one Lyapunov time from 0.34 to 0.96.
    The mechanism is anchoring: the decoder's linear interpolation maps early output indices to decorrelated input positions.

  \item \textbf{At five Lyapunov times, all methods collapse.}
    No architecture extracted predictable structure beyond this horizon on the Lorenz system.
    This is consistent with the theoretical limit imposed by sensitive dependence on initial conditions.

  \item \textbf{Degree-2 univariate NGRC underperforms.}
    The recursive NGRC ridge model reached $R^2 = 0.693$ at 0.23 Lyapunov times but only 0.157 at one Lyapunov time, consistent with the documented insufficiency of degree-2 monomials under partial observation.
\end{enumerate}


% ════════════════════════════════════════════════════════════════
\section{Limitations and Next Steps}

\begin{enumerate}[leftmargin=*]
  \item \textbf{Coarse horizon sweep.} Only three horizons were run; the crossing point of $R^2 = 0$ cannot be localised. Intermediate horizons (55, 220) and longer ones (1100, 2200) are required.
  \item \textbf{Underpowered long horizon.} $H = 550$ had 27 validation and 54 test windows and did not converge. That row should not be treated as a measurement.
  \item \textbf{Fixed error floor.} $R^2$ is flat across $H = 1, 7, 10$, indicating a decoder-side resolution limit independent of horizon.
  \item \textbf{Interpolation decoder.} The identified anchoring pathology was mitigated by shortening the input, not removed. A learned projection from the encoder's final state would address it at the source.
  \item \textbf{Non-causal filtering before splitting.} \texttt{filtfilt} and a globally estimated SWT threshold allow mild information flow across split boundaries.
  \item \textbf{Univariate NGRC.} Degree-2 monomials are known to be insufficient under partial observation; a Chebyshev basis is the documented remedy.
  \item \textbf{Additional chaotic systems.} Validation on R\"ossler, Mackey--Glass, and Lorenz-like systems would strengthen generality claims.
\end{enumerate}


% ════════════════════════════════════════════════════════════════
\begin{thebibliography}{99}

\bibitem{dudukcu2023}
H.~V. Dudukcu, M.~Taskiran, Z.~G.~C. Taskiran, and T.~Yildirim, ``Temporal Convolutional Networks with RNN approach for chaotic time series prediction,'' \emph{Applied Soft Computing}, vol.~133, p.~109945, 2023.

\bibitem{fu2024}
K.~Fu, H.~Li, and X.~Shi, ``An encoder--decoder architecture with Fourier attention for chaotic time series multi-step prediction,'' \emph{Applied Soft Computing}, vol.~156, p.~111409, 2024.

\bibitem{gauthier2021}
D.~J. Gauthier, E.~Bollt, A.~Griffith, and W.~A.~S. Barbosa, ``Next generation reservoir computing,'' \emph{Nature Communications}, vol.~12, p.~5564, 2021.

\bibitem{phien2023}
H.~N. Phien, T.~T. Dang, and N.~Q.~K. Le, ``Strategies of multi-step-ahead forecasting for chaotic time series using autoencoder and LSTM neural networks: A comparative study,'' in \emph{Proc.\ IPMV}, 2023.

\bibitem{ramadevi2022}
B.~Ramadevi and K.~Bingi, ``Chaotic time series forecasting approaches using machine learning techniques: A review,'' \emph{Symmetry}, vol.~14, no.~5, p.~955, 2022.

\bibitem{shahi2022}
S.~Shahi, F.~H. Fenton, and E.~M. Cherry, ``Prediction of chaotic time series using recurrent neural networks and reservoir computing techniques: A comparative study,'' \emph{Machine Learning with Applications}, vol.~8, p.~100300, 2022.

\end{thebibliography}

\end{document}
